\section*{Chapter \ref{ch:s2:quantitative-investment-strategies} --- Quantitative Investment Strategies}

\begin{solution}{exo:s2:quantitative-investment-strategies:1}
$0.007 / 0.10 = 0.07$ of Sharpe ratio.
\end{solution}

\begin{solution}{exo:s2:quantitative-investment-strategies:2}
$1/\sqrt{10} \approx 0.32$: a backtest Sharpe ratio of 0.5 is only about 1.6 standard errors from zero before any selection.
\end{solution}

\begin{solution}{exo:s2:quantitative-investment-strategies:3}
$1 - 0.16/0.51 = 68.6\%$ on the means; the median of each index's own decay is 72.2\%, because the ratio is skewed by indices with small backtests.
\end{solution}

\begin{solution}{exo:s2:quantitative-investment-strategies:4}
Live returns were not used to choose the index, so their error averages to zero; the backtest was chosen for being the highest of twenty, so it
carries the largest positive errors. Selection biases the statistic it selects on, not the future.
\end{solution}

\begin{solution}{exo:s2:quantitative-investment-strategies:5}
Each rule, filter and parameter added is another set of trials, and the complex strategy's backtest is the best of more of them; its errors are
larger and more of its backtest is noise. Complexity may also fit features of one period that do not recur.
\end{solution}

\begin{solution}{exo:s2:quantitative-investment-strategies:6}
Anyone who knows when the index rebalances and roughly what it will trade can buy before it buys and sell before it sells; the index then trades at
worse prices than a backtest at historical closes assumes.
\end{solution}

\begin{solution}{exo:s2:quantitative-investment-strategies:7}
The best backtest rises from 0.26 with 2 candidates to 0.76 with 100; the live result after costs from 0.07 to 0.24; the gap from 0.18 to 0.52. The
median decay stays between 68\% and 76\% because more candidates also find better strategies: the live result rises too, if much less than the
backtest. The gap in Sharpe ratio, not the percentage, is what complexity inflates here.
\end{solution}

\begin{solution}{exo:s2:quantitative-investment-strategies:8}
Passing says the best backtest is unlikely to be pure luck, not that it is accurate: those that passed backtested at 1.13 and lived at 0.32 before
costs. The selection bias remains in the backtest, and fees and leakage come on top.
\end{solution}

\begin{solution}{pb:s2:quantitative-investment-strategies:1}
\begin{enumerate}
\item A bank's rules-based strategy published as an index and accessed through swaps, notes or funds; the wrapper is the form through which a client
holds it.
\item Fees, the margin on swaps and notes, and netting the index's hedges with its other flows.
\item 2\,000 teams, 20 candidates each with true Sharpe ratios around 0.10 (sd 0.15) and errors correlated at 0.5, ten-year backtests, the best three
launched and lived for five years at 10\% volatility.
\item The distribution of true Sharpe ratios, so the median decay is near the 73\% reported for bank-built strategies.
\item 0.07 of Sharpe ratio for 0.7\% a year at 10\% volatility.
\item Nobody paid them before the index existed, and a backtest trades at historical prices nobody front-ran.
\item By running the index's positions against the swap.
\item Others trade ahead of the published rebalancing; the index buys higher and sells lower.
\item The fall from the backtest Sharpe ratio to the live after-cost Sharpe ratio across launched indices.
\item Best three: backtest 0.51, true 0.22, live 0.23 before costs and 0.16 after; median decay 72.2\%. Best one: 0.62, 0.26, 0.27, 0.20, 69.4\%.
\item 36.6\% have a negative live Sharpe ratio after costs.
\item It passes 5.9\% of teams' best indices; those backtested at 1.13 and lived at 0.32, the others at 0.59 and 0.27 before costs.
\item The launched indices backtest at a Sharpe ratio of 0.51 and live at 0.16 after costs: a gap of 0.35 and a median decay of 72\%.
\item A median 73\% fall in Sharpe ratio from backtest to live across 215 bank strategies, over 30 points more for the most complex.
\item More candidates raise the best backtest from 0.26 to 0.76 and the live result from 0.07 to 0.24: the gap grows from 0.18 to 0.52.
\item Published predictors' returns were 26\% lower out of sample and 58\% lower after publication.
\item As the best of an unknown number of trials, before fees: ask how many were tried, deflate, deduct costs and look for live years.
\item The volatility-carry index: equity volatility exposures were among the more robust live.
\item This chapter is Book~7's overfitting, measured on products that were sold.
\item Access to a strategy, packaged, costed and hedged, sold on a backtest.
\end{enumerate}
\end{solution}

\begin{solution}{iq:s2:quantitative-investment-strategies:1}
A rules-based strategy published as an index by a bank, with a rulebook and a calculation agent. The bank that pays the index's return on a swap hedges
by holding the index's positions, rebalancing them as the rules say.
\end{solution}

\begin{solution}{iq:s2:quantitative-investment-strategies:2}
How many variants were tried and how the rules were chosen; whether costs, fees and realistic execution are included; what the deflated Sharpe
ratio is; how it did in sub-periods and out of sample; and whether any live record exists.
\end{solution}

\begin{solution}{iq:s2:quantitative-investment-strategies:3}
Spread the rebalance over several days, randomise its timing within a window, trade on fixings less predictable than the close, or net it
internally; each must be written into the rulebook before it applies.
\end{solution}

\begin{solution}{iq:s2:quantitative-investment-strategies:4}
Hedging slippage against the index level, gap and jump risk on rebalancing days, counterparty risk on the swaps, model risk in complex indices,
reputational risk from poor live performance, and conflicts between the index sponsor and its trading desk.
\end{solution}

\begin{solution}{iq:s2:quantitative-investment-strategies:5}
That the index level follows the rulebook exactly, from auditable input prices; that rebalances and corrections are published as the rules say;
that historical levels are never silently restated; and that the calculation is independent of the trading desk.
\end{solution}

\begin{solution}{iq:s2:quantitative-investment-strategies:6}
The best of $N$ independent normal errors has expectation about $s\sqrt{2\ln N}$ for large $N$ (the extreme-value approximation); the live estimate
of the chosen candidate has expectation equal to its truth, the same for all. The backtest therefore exceeds the live expectation by about
$s\sqrt{2\ln N}$: 0.77 for $N = 20$ and $s = 1/\sqrt{10}$, before the slower-converging correction.
\end{solution}
